\chapter{Target Reasoning with Vision-Language Models}
\chaptermark{Target Reasoning}
\label{chap:target-reasoning}
\minitoc

\begin{center}
\begin{minipage}[b]{0.9\linewidth}
\small
\textbf{Source note.}
This chapter will be adapted from the ongoing Qwen / target reasoning manuscript. Figures and final results are expected to evolve, so this chapter is intentionally scaffolded more lightly than the previous mature paper chapters.
\end{minipage}
\end{center}

\section{Chapter Overview}
\label{sec:target-reasoning-overview}

\textbf{Drafting goal:} frame target prediction as the persistent bottleneck in surgical action triplet understanding.

\section{Why Target Prediction Is Difficult}
\label{sec:target-reasoning-difficulty}

\textbf{Drafting goal:} discuss anatomical ambiguity, visual similarity, occlusion, context dependence, and the limits of local visual representation.

\section{From Target Fusion to Target Reasoning}
\label{sec:target-reasoning-hypothesis}

\textbf{Drafting goal:} articulate the hypothesis that target prediction may require explicit reasoning over anatomy, instrument position, procedural context, and surgical intent.

\section{Reasoning-Based Framework}
\label{sec:target-reasoning-framework}

\textbf{Drafting goal:} describe the pipeline: instrument instances, multimodal model input, weak anatomical priors, structured prompting, and final triplet prediction.

\section{Target-Centric Chain-of-Thought}
\label{sec:target-reasoning-tccot}

\textbf{Drafting goal:} define the target-centric reasoning procedure, including evidence gathering, ambiguity handling, alternative hypotheses, and final prediction.

\section{Experiments}
\label{sec:target-reasoning-experiments}

\textbf{Drafting goal:} leave space for final datasets, baselines, prompts, metrics, and result tables.

\section{Discussion}
\label{sec:target-reasoning-discussion}

\textbf{Drafting goal:} discuss what reasoning models reveal about surgical target prediction and how this reframes the earlier TargetFusionNet findings.
